\chapter{Transformada euclidiana e inversión}\label{cap:v4-07}
La transformada en la recta reemplaza las frecuencias enteras por un parámetro continuo. Conservaremos el signo positivo del curso original y mantendremos una notación distinta para la transformada unitaria que aparecerá en el capítulo siguiente. Las pruebas se desarrollan en $\R^d$ cuando no hay una dificultad adicional; los argumentos específicos de la recta se identifican expresamente.

\section{La transformada de funciones integrables}\label{lectura:v4-10}
\index{Fourier!transformada euclidiana}\index{Riemann--Lebesgue, lema de}
Para $f\in L^1(\R^d)$ definimos
\[
 \widehat f(\xi)=\int_{\R^d}f(x)e^{ix\cdot\xi}\,dx.
\]
La integral existe para todo $\xi$, $|\widehat f(\xi)|\leq\norm f_1$ y no depende del representante de $f$. El signo difiere del empleado para los coeficientes periódicos; cambiarlo únicamente refleja la variable de frecuencia.
\begin{teorema}\label{v4:fourier-L1}
La transformada es una aplicación lineal continua de norma uno de $L^1(\R^d)$ en $C_0(\R^d)$. Además
\[
 \widehat{f*g}=\widehat f\,\widehat g,\quad
 \widehat{f(\cdot+a)}(\xi)=e^{-ia\cdot\xi}\widehat f(\xi),\quad
 \widehat{e^{ia\cdot x}f(x)}(\xi)=\widehat f(\xi+a).
\]
\end{teorema}
\begin{proof}
La linealidad y la cota de norma proceden de la integral. Si $f\geq0$, $\widehat f(0)=\norm f_1$, de modo que la norma del operador es uno. Para continuidad uniforme,
\[
 \sup_\xi|\widehat f(\xi+h)-\widehat f(\xi)|
 \leq\int|f(x)|\,|e^{ix\cdot h}-1|\,dx\longrightarrow0
\]
por convergencia dominada, con mayorante $2|f|$.

Para la indicadora de un rectángulo acotado $Q=\prod_j[a_j,b_j]$, Fubini expresa la transformada como el producto de las integrales unidimensionales $\int_{a_j}^{b_j}e^{ix_j\xi_j}dx_j$. Cada factor está acotado por $b_j-a_j$ y, si $\xi_j\ne0$, por $2/|\xi_j|$. Si $|\xi|\to\infty$, alguna coordenada tiene módulo al menos $|\xi|/\sqrt d$, y el producto tiende a cero. Las funciones escalonadas en rectángulos son densas en $L^1(\R^d)$ por regularidad de Lebesgue; la cota $\norm{\widehat f-\widehat g}_\infty\leq\norm{f-g}_1$ extiende la desaparición en infinito a todo $f$.

Para convolución, la integral doble de $|f(y)g(x-y)|$ es $\norm f_1\norm g_1$. Fubini y $z=x-y$ dan
\[
 \widehat{f*g}(\xi)=\int\int f(y)g(z)e^{i(y+z)\cdot\xi}\,dzdy
 =\widehat f(\xi)\widehat g(\xi).
\]
Los cambios de variable $y=x+a$ y la identidad $e^{ia\cdot x}e^{i\xi\cdot x}=e^{i(\xi+a)\cdot x}$ prueban las dos últimas fórmulas.
\end{proof}

\section{Caracteres y derivación}\label{lectura:v4-11}
\index{caracter@carácter!de la recta}
\begin{proposicion}\label{v4:caracteres-R}
Todo homomorfismo continuo $\chi:(\R,+)\to\C^\times$ tiene la forma $\chi(t)=e^{zt}$, $z\in\C$. Si toma valores en $\T$, entonces $z=i\xi$ para un único $\xi\in\R$. Los caracteres de $\R^d$ son exactamente $x\mapsto e^{ix\cdot\xi}$.
\end{proposicion}
\begin{proof}
Sea $F(t)=\int_0^t\chi(s)ds$. Como $F'(0)=\chi(0)=1$, existe $a$ con $F(a)\ne0$. La propiedad de homomorfismo da
\[
 F(a)\chi(t)=\int_t^{t+a}\chi(s)ds=F(t+a)-F(t).
\]
Así $\chi$ es diferenciable. Dividir $\chi(t+h)-\chi(t)=\chi(t)(\chi(h)-1)$ entre $h$ y pasar al límite da $\chi'(t)=\chi'(0)\chi(t)$. La derivada de $e^{-zt}\chi(t)$ es cero para $z=\chi'(0)$; su valor en cero es uno, por lo que $\chi(t)=e^{zt}$. Si $|\chi(t)|=1$, entonces $e^{t\Re z}=1$ para todo $t$, de donde $\Re z=0$. La derivada en cero prueba unicidad. En $\R^d$, restringimos a cada eje y multiplicamos las imágenes de $x_1e_1+\cdots+x_de_d$.
\end{proof}
\begin{proposicion}\label{v4:derivacion-fourier}
Si $x_jf\in L^1(\R^d)$, entonces $\partial_{\xi_j}\widehat f=\widehat{ix_jf}$. En la recta, si $f$ es localmente absolutamente continua y $f,f'\in L^1$, entonces
\[
 \widehat{f'}(\xi)=-i\xi\widehat f(\xi),\qquad
 f(x)=-\int_x^\infty f'(t)dt=\int_{-\infty}^xf'(t)dt.
\]
\end{proposicion}
\begin{proof}
El cociente incremental del exponencial está acotado por $|x_j|$, pues $|e^{iu}-1|\leq|u|$. Convergencia dominada permite derivar bajo la integral; además el error del cociente incremental está dominado independientemente de $\xi$, de modo que la convergencia es uniforme en la frecuencia.

Para la segunda afirmación, la continuidad absoluta local da $f(b)-f(a)=\int_a^bf'$. Como $f'$ es integrable, los límites de $f$ en ambos infinitos existen. Deben ser cero: un límite no nulo contradice $f\in L^1$. Pasar al límite en la identidad anterior da las dos fórmulas. La integración por partes en $[-R,R]$ produce
\[
 \int_{-R}^Rf'(x)e^{ix\xi}dx
 =[f(x)e^{ix\xi}]_{-R}^R-i\xi\int_{-R}^Rf(x)e^{ix\xi}dx.
\]
El término de frontera tiende a cero y las integrales convergen absolutamente. La hipótesis de continuidad absoluta local se hace explícita porque la mera existencia de una derivada integrable no justifica automáticamente el teorema fundamental de Lebesgue.
\end{proof}

\section{La gaussiana como puente de inversión}
\index{gaussiana}\index{nucleo@núcleo!del calor}
\begin{lema}\label{v4:gaussiana}
Se cumple $\int_\R e^{-x^2}dx=\sqrt\pi$ y
\[
 \widehat{e^{-a|x|^2}}(\xi)=\left(\frac\pi a\right)^{d/2}e^{-|\xi|^2/(4a)}\quad(a>0).
\]
Por tanto, para $t>0$,
\[
 H_t(x)=(4\pi t)^{-d/2}e^{-|x|^2/(4t)},\quad
 \widehat H_t(\xi)=e^{-t|\xi|^2},\quad
 H_t(x)=(2\pi)^{-d}\int e^{-t|\xi|^2}e^{-ix\cdot\xi}d\xi.
\]
\end{lema}
\begin{proof}
La integral $I=\int e^{-x^2}dx$ es positiva y finita por comparación en las colas con $e^{-|x|}$. Tonelli y coordenadas polares dan
\[
 I^2=\int_{\R^2}e^{-(x^2+y^2)}dxdy
 =2\pi\int_0^\infty re^{-r^2}dr=\pi.
\]
Si $F(\xi)=\int e^{-x^2}e^{ix\xi}dx$, la proposición anterior y $(e^{-x^2})'=-2xe^{-x^2}$ dan $F'(\xi)=-\xi F(\xi)/2$. Puesto que $F(0)=\sqrt\pi$, se obtiene $F(\xi)=\sqrt\pi e^{-\xi^2/4}$. El cambio $u=\sqrt a\,x$ y la factorización en coordenadas por Fubini prueban la fórmula general. Aplicar esa fórmula directamente a la gaussiana en la variable $\xi$ da la última identidad; todavía no se ha usado ningún teorema de inversión.
\end{proof}
Las funciones $H_t$ son no negativas y tienen integral uno. Para cada $\delta>0$, el cambio $x=\sqrt t\,u$ muestra $\int_{|x|\geq\delta}H_t(x)dx\to0$ cuando $t\downarrow0$. Forman, por tanto, una identidad aproximada y $f*H_t\to f$ en $L^p$ para $p<\infty$.
\begin{teorema}[Inversión]\label{v4:inversion-Rn}
Si $f\in L^1(\R^d)$ y $\widehat f\in L^1(\R^d)$, entonces
\[
 f(x)=(2\pi)^{-d}\int_{\R^d}\widehat f(\xi)e^{-ix\cdot\xi}d\xi
\]
casi en todas partes. La igualdad vale en todo punto de continuidad de $f$ y, más generalmente, en todo punto de Lebesgue.
\end{teorema}
\begin{proof}
La fórmula explícita de $H_t$ y Fubini, cuya hipótesis se verifica con $\norm f_1\int e^{-t|\xi|^2}d\xi<\infty$, dan
\[
 f*H_t(x)=(2\pi)^{-d}\int\widehat f(\xi)e^{-t|\xi|^2}e^{-ix\cdot\xi}d\xi.
\]
Si $\widehat f\in L^1$, convergencia dominada hace tender el lado derecho uniformemente en $x$ a la integral de inversión: la cota del error es $(2\pi)^{-d}\int|\widehat f||e^{-t|\xi|^2}-1|$. Por otro lado, $f*H_t\to f$ en $L^1$. Una sucesión $t_j\downarrow0$ posee una subsucesión convergente casi en todas partes, pues se puede escoger con errores $L^1$ sumables y aplicar Tonelli. La unicidad del límite prueba la fórmula casi en todas partes.

Precisemos la afirmación puntual. Sea $x$ punto de Lebesgue y $A(r)=\int_{|y|<r}|f(x-y)-f(x)|dy=o(r^d)$. Escribamos $h_t(r)=(4\pi t)^{-d/2}e^{-r^2/(4t)}$. Como $h_t(r)=\int_r^\infty-h_t'(s)ds$, Tonelli da
\[
 \int |f(x-y)-f(x)|H_t(y)dy
 =\int_0^\infty[-h_t'(s)]A(s)ds.
\]
Dado $\eps>0$, elegimos $\delta$ con $A(s)\leq\eps v_ds^d$ para $s<\delta$, donde $v_d$ es el volumen de la bola unidad. La integral en ese intervalo es a lo sumo $\eps$, porque $\int_0^\infty[-h_t'(s)]v_ds^d ds=\int H_t=1$. Para $s\geq\delta$ usamos $A(s)\leq\norm f_1+|f(x)|v_ds^d$. La contribución queda acotada por
\[
 \norm f_1h_t(\delta)+|f(x)|\left(v_d\delta^dh_t(\delta)+\int_{|y|\geq\delta}H_t(y)dy\right),
\]
integrando por partes el segundo término. Esta cota tiende a cero: la exponencial $e^{-\delta^2/(4t)}$ domina cualquier potencia de $1/t$ y la cola gaussiana tiende a cero. Así $f*H_t(x)\to f(x)$ y la fórmula puntual queda probada. Todo punto de continuidad es de Lebesgue.
\end{proof}
La prueba mediante calor sigue el enfoque de Maximenko~\cite{maximenkoInversion}, desarrollando también las constantes y el argumento en $\R^d$. En particular se aplica al error $f(x-y)-f(x)$, indispensable en el cálculo del punto de Lebesgue.

\section{Inversión truncada y localización en la recta}
\begin{proposicion}[Criterio de Dini en $\R$]\label{v4:dini-R}
Sea $f\in L^1(\R)$ y $S_Rf(x)=(2\pi)^{-1}\int_{-R}^R\widehat f(\xi)e^{-ix\xi}d\xi$. Si
\[
 \int_0^\delta\frac{|f(x+t)+f(x-t)-2f(x)|}{t}dt<\infty
\]
para algún $\delta>0$, entonces $S_Rf(x)\to f(x)$.
\end{proposicion}
\begin{proof}
Fubini se aplica al dominio finito de frecuencias, pues la integral absoluta es a lo sumo $2R\norm f_1$. Integrando el exponencial se obtiene
\[
 S_Rf(x)=\frac1\pi\int_0^\infty[f(x+t)+f(x-t)]\frac{\sin Rt}{t}dt.
\]
Usamos la integral impropia $\int_0^\infty\sin(Rt)/t\,dt=\pi/2$ para $R>0$. Para verificarla, la integral amortiguada $J_a(R)=\int_0^\infty e^{-at}\sin(Rt)/t\,dt$ satisface $J_a'(R)=a/(a^2+R^2)$ y $J_a(0)=0$, luego $J_a(R)=\arctan(R/a)$. Integración por partes da la cota uniforme $|\int_M^Ne^{-at}\sin(Rt)/t\,dt|\leq3/(RM)$ para $a\geq0$. Se pasa primero $a\downarrow0$ en un intervalo finito y después $M\to\infty$, obteniendo $\pi/2$.

Restamos $f(x)$ usando esa integral y separamos $(0,\delta)$ del resto. En el primer intervalo, la condición de Dini y Riemann--Lebesgue hacen tender a cero la integral del error. En $(\delta,\infty)$, $[f(x+t)+f(x-t)]/t$ es integrable, por lo que su integral oscilante tiende a cero. El término constante restante es un múltiplo de $\int_\delta^\infty\sin(Rt)/t\,dt$, acotado por $3/(R\delta)$ y por ello convergente a cero. Es necesario tratar este término aparte: la función constante dividida entre $t$ no es integrable en infinito, contrario al uso directo de Riemann--Lebesgue sugerido en la lectura 10.
\end{proof}

\section{Ejercicios de la transformada integrable}
\begin{ejercicio}[Lectura 10, ejercicio 1: transferencia]\label{ej:v4-l10-1}
Si $f,g\in L^1(\R^d)$, probar $\int f\widehat g=\int\widehat f g$.
\end{ejercicio}

\begin{ejercicio}[Lectura 10, ejercicio 2: involución]\label{ej:v4-l10-2}
Definir $f^*(x)=\overline{f(-x)}$. Probar que es una involución isométrica de $L^1(\R^d)$ y que $\widehat{f^*}=\overline{\widehat f}$. Explicar por qué $L^1(\R)$ no es un álgebra $C^*$ con esa norma.
\end{ejercicio}

\begin{ejercicio}[Lectura 10, ejercicio 3]\label{ej:v4-l10-3}
Si $(k_j)$ es una sucesión de Dirac en $L^1(\R^d)$, demostrar $\widehat k_j(\xi)\to1$ y que, para todo $h\in C_0(\R^d)$, $\norm{h\widehat k_j-h}_\infty\to0$.
\end{ejercicio}


\section{Ejercicios de inversión y funciones de Schwartz}
\index{Schwartz!espacio de}
\begin{ejercicio}[Lectura 11, ejercicio 1]\label{ej:v4-l11-1}
Si $f\in L^1(\R^d)$ y $\widehat f=0$, demostrar que $f=0$ casi en todas partes.
\end{ejercicio}

\begin{ejercicio}[Lectura 11, ejercicio 2: factor de normalización]\label{ej:v4-l11-2}
Calcular la transformada de $\mathbf1_{[-1,1]}$ y demostrar que no es integrable.
\end{ejercicio}

\begin{ejercicio}[Lectura 11, ejercicio 3: Parseval integrable]\label{ej:v4-l11-3}
Si $f,g,\widehat f,\widehat g\in L^1(\R^d)$, demostrar
$\int\widehat f\,\overline{\widehat g}=(2\pi)^d\int f\overline g$.
\end{ejercicio}

\begin{ejercicio}[Lectura 11, ejercicio 4]\label{ej:v4-l11-4}
Para $\varphi(x)=e^{-x^2}$, calcular su integral y transformada. Probar que $k_n(x)=n\pi^{-1/2}e^{-n^2x^2}$ forma una sucesión de Dirac que satisface la fórmula de inversión.
\end{ejercicio}

\begin{ejercicio}[Lectura 11, ejercicio 5]\label{ej:v4-l11-5}
Una función $f\in C^\infty(\R^d)$ es de Schwartz si $\sup_x|x^\alpha\partial^\beta f(x)|<\infty$ para todos los multiíndices $\alpha,\beta$. Probar: (i) $C_c^\infty\subset\Sscr$ y las gaussianas pertenecen a $\Sscr$ sin tener soporte compacto; $\Sscr$ es denso en $L^p$, $p<\infty$; (ii) Fourier preserva $\Sscr$; (iii) $L^1(\R)$ tiene divisores de cero.
\end{ejercicio}


\section{Los caracteres del álgebra de convolución}\label{lectura:v4-12}
\index{caracter@carácter!de L1}\index{Gelfand!transformada de}
Un carácter de un álgebra de Banach compleja es un funcional lineal multiplicativo, continuo y no nulo. Cada evaluación $\phi_\xi(f)=\widehat f(\xi)$ es un carácter de $L^1(\R^d)$ por el teorema de convolución; es no nulo, por ejemplo, sobre una gaussiana modulada.
\begin{lema}
Todo carácter $\phi$ de un álgebra de Banach satisface $\norm\phi\leq1$.
\end{lema}
\begin{proof}
Si $|\phi(a)|>\norm a$, pongamos $b=a/\phi(a)$. Entonces $\phi(b)=1$ y $\norm b<1$. La serie $x=-\sum_{n\geq1}b^n$ converge absolutamente y satisface $b+x=bx$. Aplicar $\phi$ daría $1+\phi(x)=\phi(x)$, imposible. El argumento no requiere una unidad en el álgebra.
\end{proof}
\begin{teorema}\label{v4:caracteres-L1R}
Todo carácter de $L^1(\R^d)$ es $f\mapsto\widehat f(\xi)$ para una única frecuencia $\xi$. Esta identificación es un homeomorfismo entre $\R^d$ y el espacio de caracteres con su topología débil estrella.
\end{teorema}
\begin{proof}
Por la representación del dual de $L^1(\R^d)$ existe $h\in L^\infty$ con $\phi(f)=\int fh$. Para $R_yf(x)=f(x-y)$, Fubini y multiplicatividad dan, para cada $f,g\in L^1$,
\[
 \int g(y)\phi(R_yf)dy=\phi(f*g)=\phi(f)\phi(g)=\int g(y)\phi(f)h(y)dy.
\]
La integral doble es absolutamente convergente por $\norm h_\infty\norm f_1\norm g_1$. Como $g$ es arbitraria, $\phi(R_yf)=\phi(f)h(y)$ para casi todo $y$.

Elijamos $f_0$ con $\phi(f_0)\ne0$ y definamos $H(y)=\phi(R_yf_0)/\phi(f_0)$. La continuidad de traslaciones hace a $H$ continua y acotada, y $H=h$ casi en todas partes. Para cualquier $f$, las funciones continuas $y\mapsto\phi(R_yf)$ y $y\mapsto\phi(f)H(y)$ coinciden casi en todas partes, luego en todo punto: una diferencia continua no nula sería no nula en un abierto de medida positiva. Así se obtiene una identidad válida para cada $f$ y todo $y$, sin elegir ilegítimamente un único nulo para todos los $f$.

Aplicándola a $R_zf_0$, se deduce $H(y+z)=H(y)H(z)$ y $H(0)=1$. En particular $H$ nunca se anula. La acotación y las potencias positivas y negativas implican $|H(y)|=1$; si fuese mayor o menor que uno, alguna sucesión de potencias sería no acotada. Por la proposición~\ref{v4:caracteres-R}, $H(y)=e^{iy\cdot\xi}$, y por tanto $\phi=\phi_\xi$. Si dos frecuencias dan el mismo funcional, sus exponenciales son iguales casi en todas partes por dualidad y luego en todo punto por continuidad; restringir a los ejes y derivar en cero prueba que las frecuencias coinciden.

La continuidad de $\xi\mapsto\phi_\xi$ es la continuidad de cada $\widehat f$. Para la inversa, si una red $\phi_{\xi_\alpha}\to\phi_\xi$ débil estrella, evaluamos en $f(x)=e^{-|x|^2}e^{-ix\cdot\xi}$. Su transformada es $\pi^{d/2}e^{-|\eta-\xi|^2/4}$. La convergencia de esas evaluaciones implica $e^{-|\xi_\alpha-\xi|^2/4}\to1$, y de ahí $\xi_\alpha\to\xi$. Esto prueba continuidad inversa incluso para redes.
\end{proof}
La transformada de Gelfand asigna a $f$ la función $\phi\mapsto\phi(f)$ sobre el espacio de caracteres. Bajo el homeomorfismo anterior es exactamente la transformada de Fourier. La lectura 12 no contiene una lista de ejercicios; sus resultados de norma, continuidad del representante y topología del espectro se han demostrado aquí.

\paragraph{Del álgebra a la geometría.}
La transformada ya distingue funciones integrables y respeta su producto de convolución. El siguiente paso cambia de norma: Plancherel extenderá Fourier a todo $L^2$, incluso cuando la integral que lo define no converja absolutamente.
